\section{量子计算基本理论}
量子计算利用了量子力学的两个主要特性：叠加与纠缠。

\subsection{量子叠加}
量子叠加是指一个量子系统可以同时存在于多个状态或位置。这是量子计算并行处理能力的基础。

单个量子位的状态可以表示为二维向量空间 $\mathbb{C}^2$ 中的单位列向量：
\[
\ket{\psi} = \alpha\ket{0} + \beta\ket{1}
\]
其中 $\alpha,\beta \in \mathbb{C}$，且满足 $|\alpha|^2 + |\beta|^2 = 1$。

\subsection{量子纠缠}
量子纠缠是指两个或多个量子粒子之间存在的一种极强的关联。无论这些粒子相距多远，它们的状态都能瞬间完全一致地改变，多个量子位组成的联合状态需要用考虑了量子位间相互作用（即纠缠）的新向量空间来描述。

\subsection{张量积空间}
多量子位量子系统可以由多个已知的量子子系统构成。设第 $i$ 个子系统的状态空间为可分离的希尔伯特空间 $H_i$，那么多量子位系统的状态空间就是张量积空间：
\[
H = H_1 \otimes H_2 \otimes \cdots \otimes H_n
\]

多量子位系统的基态由单量子位的基态通过向量的张量积构成：
\[
\ket{\psi_1} \otimes \ket{\psi_2} \otimes \cdots \otimes \ket{\psi_n}
\]

\subsection{量子寄存器}
量子寄存器由若干量子位组成，寄存器的大小由量子位的数量决定。$m$ 个量子位可以同时表示 $2^m$ 个状态。

例如，4量子位寄存器可以由16种状态的叠加来表示：
\[
\ket{\psi} = c_0\ket{0000} + c_1\ket{0001} + \cdots + c_{15}\ket{1111}
\]
